Time series analysis is one of the oldest and most economically consequential branches of statistical machine learning. Every electricity grid, ICU, financial exchange, semiconductor fab, satellite constellation, and supply chain produces high-frequency real-valued measurements. These measurements must be summarised, forecast, monitored for anomalies, classified, or imputed. For decades, the dominant tools were parametric statistical models --- autoregressive integrated moving average (ARIMA), exponential smoothing (ETS), and state-space models such as the Kalman filter --- combined with shallow tree- and kernel-based learners. The deep-learning era introduced recurrent networks (LSTMs, GRUs) and, after 2017, convolutional and Transformer-based architectures that pushed long-horizon accuracy on standard benchmarks such as ETTh1, Weather, Electricity, and Traffic. Between 2021 and 2023 the community produced a rapid succession of specialised time-series Transformers --- Informer (Zhou et al., AAAI 2021) {[}1{]}, Autoformer (Wu et al., NeurIPS 2021) {[}2{]}, FEDformer (Zhou et al., ICML 2022) {[}3{]}, PatchTST (Nie et al., ICLR 2023) {[}4{]}, and iTransformer (Liu et al., ICLR 2024) {[}5{]} --- yet a striking parallel observation by Zeng et al.~in their AAAI 2023 paper ``Are Transformers Effective for Time Series Forecasting?'' showed that a single linear layer with the right normalisation could match or beat them on six widely-used long-term forecasting benchmarks {[}6{]}. Against this backdrop of architect\ldots{}
